\label{sec:other}

There are a bunch of other scenarios that take the subjunctive that don't quite clearly fit into the WEIRDO categories. In each of the following sections, remember to think about why the subjunctive is being used (typically to add an element of uncertainty).

\section{\ita{Si}-Clauses}

What's better than Santa Claus? A \ita{si}-clause! These are pretty fun, so let's start with them. \\

\ita{Si}-clauses are used to expresse if-then scenarios. There are exactly two different options for us here:
\begin{itemize}[noitemsep]
	\item The \ita{si}-clause is in the present indicative, in which case the subordinate clause will be in the present indicative.
	\item The \ita{si}-clause is in the subjunctive imperfect, in which case the subordinate clause will be in the conditional. 
\end{itemize}

I'm assuming you've already worked with sentences in the present indicative, so let's look at some sample sentences where we're using the subjunctive imperfect and the conditional:
\begin{enumerate}[noitemsep]
	\item ``If the painting were on the wall, everyone would look at it'' $\rightarrow$ \ita{Si el cuadro estuviera sobre la pared, todos lo ver\'ian.}
	\item ``If he had been a doctor, he would have been richer.'' $\rightarrow$ \ita{Si fuera doctor, habr\'ia sido m\'as rico.}
	\item ``They would be happier if the work were done.'' $\rightarrow$ \ita{Estar\'ian m\'as contentos si el trabajo estuviera terminado.}
	\item ``If I were you, I'd run far away from here.'' $\rightarrow$ \ita{Si yo fuera t\'u, correr\'ia lejos de aqu\'i.}\footnote{You can also use the phrase \ita{yo de ti} in this context, but I just wanted to show you the use of the subjunctive imperfect.}
\end{enumerate}


\section{What if?}
Another use of the subjunctive imperfect is in expressing ``what if'' scenarios. These are pretty simple: just replace the ``what if'' with \ita{y si} and conjugate the main verb in the subjunctive imperfect, and you're good to go!\\

Here are some example sentences with their translations:

\begin{enumerate}[noitemsep]
	\item ``What if it rains?'' $\rightarrow$ \ita{?`Y si lloviera?}
	\item ``What if the car gets dented?'' $\rightarrow$ \ita{?`Y si el coche se abollara?}
	\item ``What if the frog became a prince''? $\rightarrow$ \ita{?`Y si la rana se convirtiera en pr\'incipe?}
\end{enumerate}


\section{Maybe}

In Spanish, there are a couple of different ways to say ``maybe,'' \ita{quiz\'a(s)} or \ita{tal vez}. Both expressions can be followed up by either the subjunctive or the indicative, and still be considered correct. You can think of this as more or less of a spectrum -


\begin{center}
\begin{tikzpicture}

\shade[left color=yellow,right color=red] (0,0) rectangle (15,0.5);

\end{tikzpicture}
\end{center}
Certain \hfill Neutral \hfill Uncertain\\ 
 
By using the word maybe, we start out somewhere in the middle, at an orange, already. If the speaker's emotion falls on the ``uncertain'' side (reddish) of the spectrum, use the subjunctive. If the speaker's emotion falls on the certain side (yellowish) of the spectrum, use the indicative. \\

For example, the sentences \ita{Quiz\'as \'el se siente triste} and \ita{Quiz\'as \'el se sienta triste} are both grammatically correct. The former expression is used when the speaker is pretty confident that the individual is sad, and the latter is used when the speaker doesn't feel so sure.

\section{Other expressions}

The following expressions \underline{always} trigger the subjunctive. \\


\begin{table}[H]
	\centering
	\begin{tabular}{ll}
	\toprule
	\textbf{Phrase} & \textbf{Possible translation} \\
	\midrule
	so that & \ita{para que} \\
			& \ita{a fin de que} \\
			& \ita{a fin de que} \\
	%\ita{para que} & so that \\
	%\ita{a fin de que} & \\
	%\ita{con objeto de que} & \\
	without & \ita{sin que} \\ 
	%\ita{sin que} & without \\
	unless & \ita{a menos que} \\
			& \ita{a no ser que} \\ 
	%\ita{a menos que} & unless\\
	%\ita{a no ser que} & \\
	provided/given that & \ita{con tal de que} \\
						& \ita{dado que} \\
						& \ita{siempre que} \\
						& \ita{siempre y cuando} \\
						& \ita{a condici\'on de que} \\
	%\ita{con tal de que} & provided/given that \\
	%\ita{dado que} & \\
	%\ita{siempre que} & \\
	%\ita{siempre y cuando} & \\
	\bottomrule
	\end{tabular}
	\caption{A few other expressions that trigger the subjunctive}
\end{table}


Here are some sample sentences with their corresponding translations:
	\begin{enumerate}[noitemsep]
		\item ``He explained it to them so that they would understand.'' $\rightarrow $ 
			\begin{itemize}[noitemsep]
				\item \ita{Se lo explic\'o para que entendieran.}
				\item \ita{Se lo explic\'o a fin de que entendieran.}
				\item \ita{Se lo explic\'o con objeto de que entendieran.}
			\end{itemize}
		\item ``I approached the frog without\footnote{\ita{Sin que} is used when we want to use a conjugated verb after the word ``without.'' If we can follow up with an infinitive, we can just use the word \ita{sin.}} it noticing.'' $\rightarrow$ \ita{Me acerqu\'e a la rana sin que se diera cuenta.}
		\item ``He won't help you unless you pay him.'' $\rightarrow$
			\begin{itemize}[noitemsep]
				\item \ita{No te va a ayudar a menos que le pagues.}
				\item \ita{No te va a ayudar a no ser que le pagues.}
			\end{itemize}
		\item ``They'll eat anything, provided \footnote{\ita{Siempre que} takes the subjunctive when it is used to mean ``provided that.'' When it's used to mean ``whenever,'' it sometimes takes the indicative. So watch out!} it doesn't contain walnuts.'' $\rightarrow$ 
			\begin{itemize}[noitemsep]
				\item \ita{Comer\'an cualquier comida, con tal de que no contenga nueces. }
				\item \ita{Comer\'an cualquier comida, siempre que no contenga nueces.}
				\item \ita{Comer\'an cualquier comida, siempre y cuando no contenga nueces.}
				\item \ita{Comer\'an cualquier comida, a condici\'on de que no contenga nueces.}
			\end{itemize}
	\end{enumerate}

\section{Time Expressions}

A time expression uses an event to describe another event. For example, the words ``before,'' and ``after'' are typically used to set off a time expression. With these expressions, sometimes we use the subjunctive, and other times, we use the subjunctive. \\

	If you're using a time expression with an event that happened in the past, you can typically just use the indicative. For example, to say ``He washed the car before I finished my homework,'' we'd just say \ita{Lav\'o el coche antes de que termin\'e mi tarea.} To translate ``After she fell ill, he took care of all the chores,'' you would say \ita{Despu\'es de que se enferm\'o, \'el se encarg\'o de todos los quehaceres.} The reason why we use the indicative with time expressions in the past tense is that since the event has already taken place, there is no uncertainty or subjectivity involved. \\

As you might guess, when we use time expressions with events set to happen in the future, we need to use the subjunctive. This is because we can't be sure if the event will happen; no one can foresee the future. Hence, to translate the sentence ``After you finish your homework, can you turn off the lights?,'' we'd say \ita{Despu\'es de que termines tu tarea, ?`puedes apagar las luces?} This is because the speaker can't be sure that the person they're addressing will finish their homework. Similarly, to translate the sentence ``I want to do the laundry before she gets home,'' we can say \ita{Quiero hacer la colada antes de que llegue a casa.} \\

Also, just as we discussed in previous sections, we only use the subjunctive if there are two distinct subjects in the main and subordinate clauses. If we have the same subject, the infinitive works just fine. For example, to translate the sentence ``Before I forget, I'm going to shut down my computer,'' we could say \ita{Antes de olvidar, voy a apagar mi computadora/ordenador.} Most native speakers wouldn't say \ita{Antes de que olvide, voy a apagar mi computadora/ordenador.} \\

\section{Usage with other relative pronouns}

In the majority of cases, the subordinate clause will begin with the relative pronoun \ita{que}. However, the relative pronouns \ita{cuando}, \ita{donde}, and \ita{quien} can also set off a subordinate clause. In some of these instances, we may have to use the subjunctive, while in others, the indicative may be used. \\

Like before, we'll use the indicative when we feel certain, and the subjunctive when we feel uncertain. \\

Let's do an example with the relative pronoun \ita{quien.}:
\begin{enumerate}[noitemsep]
	\item ``I'm looking for someone that speaks ``Marathi'' $\rightarrow$ \ita{Busco a alguien que hable marat\'i.}
	\item ``I met a guy who speaks ``Marathi''$ \rightarrow$ \ita{Conoc\'i a un chavo que hablaba marat\'i.}
\end{enumerate}

Why are we using the subjunctive in the first sentence, but not in the second? Shouldn't we be using the indicative in the first sentence? \\

The reason we have to use the subjunctive in the first sentence is because the speaker isn't sure that there is a person (wherever they're looking) that speaks Marathi. However, in the second sentence, the speaker is certain that they met a Marathi speaker, which is why the indicative is used. \\

Let's do another example with the relative pronoun \ita{cuando}:
\begin{enumerate}
	\item ``When he returns, you can speak with him'' $\rightarrow$ \ita{Cuando llegue, puedes hablar con \'el.}
	\item ``When he returned, I spoke with him'' $\rightarrow$ \ita{Cuando lleg\'o, habl\'e con \'el.}
\end{enumerate}

Just like the previous example, the subjunctive is used when the speaker isn't that sure. In the first example, the subjunctive is used since the speaker \underline{cannot} be sure that he'll ever return (sad, I know). In the second example, since the speaker can be certain that he returned, the indicative is used. \\

Notice a pattern? In these sorts of instances, if the sentence is in the past tense, we can use the indicative in many (not all) cases, since we can generally be sure of what happened in the past. Since the present and future are more unpredictable, we generally use the subjunctive with time-related expressions like these ones. 

\section{Who/What/Which/When/Where + ever}

To convey the idea of whoever, whatever, whenever, or wherever, we use the subjunctive - ish. I'm also going to talk about the idea of ``any,'' which is a related concept, in this section. \\

\subsection{Whatever}
Let's start out with ``whatever,'' which can be translated as \ita{lo que + verbo}, where the verb is conjugated in the subjunctive mood. For example, the sentence ``You can pick whatever you want'' can be translated as \ita{Puedes escoger lo que quieras.} To translate the idea of ``whatever'' in general (when it means ``anything''), you'd use \ita{lo que sea.} For example, you translate ``I would do anything (whatever) to save frogs.'' as \ita{Har\'ia lo que sea para salvar las ranas. }

\subsection{Whichever}
\label{subsec:cualquier}
To translate ``whichever,'' we say \ita{cualquier}, the apocope of \ita{cualquiera}. There's an in depth explaination in the \hyperref[sec:apo]{Apocapation} section in the appendix. For example, to translate the sentence ``You can pick whichever frog,'' we can say \ita{Puedes elegir cualquier rana.} See how in that example, we could have said ``You can pick any frog you like'' and convey the exact same idea? This is a great time to remember that:

\begin{conf}{Language tips :)}
	Here's a general tip to help you out through your language learning journey: more often than not, there are multiple correct translations. Sometimes (just like in math), it can be useful to rephrase the problem in an equivalent form to help you think of the translation.  
\end{conf}

\begin{figure}
 \centering
 \begin{tikzpicture}[ele/.style={fill=black,circle,minimum width=.8pt,inner sep=1pt},every fit/.style={ellipse,draw,inner sep=-2pt}]

\node at (0,0) {$Language-1$};
\node at (4,0) {$Language-2$};
  \node[ele,label=left:$a$] (a1) at (0,4) {};    
  \node[ele,label=left:$b$] (a2) at (0,3) {};    
  \node[ele,label=left:$c$] (a3) at (0,2) {};
  \node[ele,label=left:$d$] (a4) at (0,1) {};

  \node[ele,,label=right:$1$] (b1) at (4,4) {};
  \node[ele,,label=right:$2$] (b2) at (4,3) {};
  \node[ele,,label=right:$3$] (b3) at (4,2) {};
  \node[ele,,label=right:$4$] (b4) at (4,1) {};

  \node[draw,fit= (a1) (a2) (a3) (a4),minimum width=2cm] {} ;
  \node[draw,fit= (b1) (b2) (b3) (b4),minimum width=2cm] {} ;  
  %\draw[->,thick,shorten <=2pt,shorten >=2pt] (a1) -- (b4);
  \draw[->,thick,shorten <=2pt,shorten >=2] (a1) -- (b2);
  \draw[->,thick,shorten <=2pt,shorten >=2] (a2) -- (b2);
  \draw[->,thick,shorten <=2pt,shorten >=2] (a3) -- (b1);
  \draw[->,thick,shorten <=2pt,shorten >=2] (a3) -- (b3);
  \draw[->,thick,shorten <=2pt,shorten >=2] (a4) -- (b4);
 \end{tikzpicture}

\caption{Visualization of translations}
\label{fig:function}
\end{figure}

\begin{conf}{Functions}
	To be honest, it's easiest to think of translations as methods, and less like functions. Most students tend to think of translatons as functions, where they're just mappings in which one input (in Language 1) maps to a unique translation in Language 2. However, this is hardly true. In the majority of cases, one input can map to several outputs, and one output can map to several inputs (as shown in figure~\ref{fig:function}). Sometimes, a translation of a particular sentence or phrase is a one-to-one mapping, but this doesn't happen as frequently as we'd hope.
\end{conf}

Use \ita{cualquier} when the meaning is somewhere between ``any'' and ``whichever.''
\subsection{Wherever}

With translating the idea of ``wherever,'' we have a few different options, where we can choose to use the subjunctive or the indicative. 

\subsubsection{With subjunctive}
With the subjunctive, ``wherever'' works just as ``whatever'' worked, with a bit of a difference. Remember how there are two different ways to ask where, \ita{d\'onde} and \ita{a d\'onde/ad\'onde}\footnote{Within this guide, I'm going to be using \ita{adonde}, not \ita{a donde}}? When we translate ``whenever,'' keep in mind whether we need to indicate movement. \\

Within this subcategories, we have a couple of different options. 
The idea of ``wherever'' can be translated as \ita{donde/adonde + verbo} or \ita{dondequiera/adondequiera que + verbo}, where the verb is conjugated in the subjunctive mood. \footnote{For the latter option, the verb in question is typically \ita{ser}.} \\

Now, our decision tree branches out a bit. In most cases, we use the \ita{donde/adonde + verbo} option. In more formal examples, such as in poetry or literature, we can use \ita{dondequiera/adondequiera que + verbo}. \\

For example, the sentence ``We can eat wherever you want'' can be translated as \ita{Podemos comer donde quieras} or \ita{Podemos comer dondequiera que sea.} The first sentence is used for most part, but the second sentence can be used in more formal contexts. Note that we're using the form related to the pronoun \ita{donde}, since we don't need to indicate movement. \\

However, if we wanted to translate the sentence ``The frogs followed the insects wherever they went,'' we could say \ita{Las ranas siguieron los insectos adonde fueran/fuesen} or \ita{Las ranas siguieron los insectos adondequiera que fueran/fuesen.} Here, since we need to indicate the frogs' movement toward the insects, we're using the \ita{adonde} form. 

\subsubsection{Without subjunctive}
There's also a small hack we can use to convey the idea of ``wherever'' without using the subjunctive: using \hyperref[subsec:cualquier]{\ita{cualquier}}. We can say \ita{cualquier lugar/sitio/parte}. 

\subsection{Whenever}


\subsection{Whoever}
Let's move on to whoever, which can be translated as \ita{quienquiera que ...}. For example, the sentence  ``Whoever you are/may be, please leave'' can be translated as \ita{Quienquiera que seas, por favor salga.} \footnote{I've chosen to use the formal command in this case, since the speaker obviously doesn't know who they're addressing.}\\

\subsubsection{Anyone}
This is also a good time to talk about the phrase ``anyone'' or ``anybody,'' 


The word \ita{cualquiera} follows \hyperref[sec:apo]{apocapation}. 

\section{Even if/though}

\ita{Aunque} can either be translated as ``even if'' or ``even though.'' \ita{Por m\'as que} can be translated as \ita{However ... it may be}. If the speaker feels pretty certain about what's about to follow, use the indicative. Else, use the subjunctive. However, you'll typically find that \ita{aunque} and \ita{por m\'as que} are \underline{typically} used with the subjunctive. \\

Let's look at a few sample sentences: 
\begin{enumerate}[noitemsep]
	\item ``Even if it's hard, don't give up.''/``However hard it may be, don't give up.'' \arr
		\begin{itemize}[noitemsep]
			\item \ita{Aunque sea dif\'icil, no te rindas.}
			\item \ita{Por m\'as dif\'icil que sea, no te rindas.}
		\end{itemize}
	\item ``Even though the clay pot broke, I enjoyed making it.'' \arr \ita{Aunque el tarro de arcilla se rompi\'o, disfrut\'e cre\'andolo.}
\end{enumerate}

\section{In case ...}
